% Generated from locale/tr/content/lambda-calculus/church-rosser/parallel-beta-reduction.tex; do not hand-edit.


\olfileid[tr]{lam}{cr}{pb}

\olsection{Paralel $\beta$-indirgeme}

\emph{Paralel $\beta$-indirgeme} kavramını tanıtıp Church--Rosser özelliğini
sağladığını ispatlayacağız.

\begin{defn}[Paralel $\beta$-indirgeme, $\bredpar$] \ollabel{defn:bredpar}
  Terimlerin paralel indirgenmesi ($\bredpar$) tümevarımlı olarak şöyle
  tanımlanır:
  \begin{enumerate}
    \item \ollabel{defn:bredpar1} $x \bredpar x$.
    \item \ollabel{defn:bredpar2} $N \bredpar N'$ ise
      $\lambd[x][N] \bredpar \lambd[x][N']$.
    \item \ollabel{defn:bredpar3} $P \bredpar P'$ ve $Q \bredpar Q'$ ise
      $PQ \bredpar P'Q'$.
    \item \ollabel{defn:bredpar4} $N \bredpar N'$ ve $Q \bredpar Q'$ ise
      $(\lambd[x][N])Q \bredpar \Subst{N'}{Q'}{x}$.
  \end{enumerate}
\end{defn}

Paralel $\beta$-indirgeme, bir terimdeki istenen sayıda redexi tek adımda
indirgememize izin verir. $\beta$-indirgemeden şu bakımdan ayrılır: yalnızca
başlangıç teriminde bulunan redexler daraltılabilir; paralel indirgeme
sonucunda ortaya çıkanlar aynı adımda daraltılamaz. Örneğin
$(\lambd[f][fx])(\lambd[y][y])$ terimi paralel olarak yalnızca kendisine ya
da $(\lambd[y][y])x$ terimine indirgenebilir; $\beta$-indirgemeyle $x$'e
indirgenebilmesine rağmen tek paralel adımda $x$'e indirgenemez. Çünkü bu
redex ancak ilk paralel $\beta$-indirgeme adımından sonra ortaya çıkar. İkinci
bir paralel adım ise $x$'i verir.

\begin{thm}\ollabel{thm:refl}
  $M \bredpar M$.
\end{thm}

\begin{proof}
  Alıştırma.
\end{proof}

\begin{prob}
  \olref[lam][cr][pb]{thm:refl} teoremini ispatlayın.
\end{prob}

\begin{defn}[$\beta$-tam geliştirme]\ollabel{defn:bcd}
  $M$ teriminin \emph{$\beta$-tam geliştirmesi} $\bcd{M}$ tümevarımlı olarak
  şöyle tanımlanır:
  \begin{align}
    \bcd{x} &= x \ollabel{defn:bcd1} \\
    \bcd{(\lambd[x][N])} &= \lambd[x][\bcd{N}] \ollabel{defn:bcd2}\\
    \bcd{(PQ)} &= \bcd{P}\bcd{Q} && \text{$P$ bir $\lambd$-soyutlaması değilse}
    \ollabel{defn:bcd3} \\
    \bcd{((\lambd[x][N])Q)} &= \Subst{\bcd{N}}{\bcd{Q}}{x} \ollabel{defn:bcd4}
  \end{align}
\end{defn}

Bir terimin $\beta$-tam geliştirmesi, adından da anlaşılacağı gibi ``tam bir
paralel indirgeme''dir. Paralel $\beta$-indirgemede bir redexi daraltmamayı
seçebiliriz; tam geliştirmede ise bütün redexleri daraltmak zorundayız.
Dolayısıyla $(\lambd[f][fx])(\lambd[y][y])$ teriminin tam geliştirmesi
kendisi değil, $(\lambd[y][y])x$'tir.

\begin{editorial}
  Bu tanımın bir sorunu vardır: ($\lambd$-)terimleri üzerinde fonksiyonların
  özyinelemeli olarak nasıl tanımlanacağını henüz açıklamadık. Bu eksiklik
  ileride giderilecektir.
\end{editorial}

\begin{lem}\ollabel{lem:comp}
  $M \bredpar M'$ ve $R \bredpar R'$ ise $\Subst{M}{R}{y}
  \bredpar \Subst{M'}{R'}{y}$.
\end{lem}

\begin{proof}
  $M \bredpar M'$ !!{derivation} üzerinde tümevarımla.
  \begin{enumerate}
    \item Son adım \olref{defn:bredpar1}: alıştırma.
    \item Son adım \olref{defn:bredpar2} olsun. $M=\lambd[x][N]$ ve
      $M'=\lambd[x][N']$; burada $N\bredpar N'$. Göstermek istediğimiz
      $\Subst{(\lambd[x][N])}{R}{y}\bredpar
      \Subst{(\lambd[x][N'])}{R'}{y}$, yani
      $\lambd[x][\Subst{N}{R}{y}]\bredpar
      \lambd[x][\Subst{N'}{R'}{y}]$ bağıntısıdır. Bu,
      \olref{defn:bredpar2} ile tümevarım varsayımından doğrudan çıkar.
    \item Son adım \olref{defn:bredpar3}: alıştırma.
    \item Son adım \olref{defn:bredpar4} olsun. Bu durumda
      $M=(\lambd[x][N])Q$ ve $M'=\Subst{N'}{Q'}{x}$. Göstermek istediğimiz
      $\Subst{((\lambd[x][N])Q)}{R}{y}\bredpar
      \Subst{\Subst{N'}{Q'}{x}}{R'}{y}$, yani
      $(\lambd[x][\Subst{N}{R}{y}])\Subst{Q}{R}{y}\bredpar
      \Subst{\Subst{N'}{R'}{y}}{\Subst{Q'}{R'}{y}}{x}$ bağıntısıdır. Bu,
      \olref{defn:bredpar4} ile tümevarım varsayımından çıkar.
  \end{enumerate}
\end{proof}

\begin{prob}
  \olref[lam][cr][pb]{lem:comp} lemasının ispatını tamamlayın.
\end{prob}

\begin{lem}\ollabel{lem:cont}
  $M \bredpar M'$ ise $M' \bredpar \bcd{M}$.
\end{lem}
\begin{proof}
  $M \bredpar M'$ !!{derivation} üzerinde tümevarımla.
  \begin{enumerate}
    \item Son kural \olref{defn:bredpar1}: alıştırma.
    \item Son kural \olref{defn:bredpar2} olsun. $M=\lambd[x][N]$ ve
      $M'=\lambd[x][N']$; burada $N\bredpar N'$. Göstermek istediğimiz
      $\lambd[x][N']\bredpar\bcd{(\lambd[x][N])}$, yani
      \olref{defn:bcd2} gereği
      $\lambd[x][N']\bredpar\lambd[x][\bcd{N}]$ bağıntısıdır. Bu,
      \olref{defn:bredpar2} ile tümevarım varsayımından çıkar.
    \item Son kural \olref{defn:bredpar3} olsun. Bazı $P,Q,P',Q'$ için
      $M=PQ$, $M'=P'Q'$, $P\bredpar P'$ ve $Q\bredpar Q'$ olur. Tümevarım
      varsayımına göre $P'\bredpar\bcd{P}$ ve $Q'\bredpar\bcd{Q}$.
      \begin{enumerate}
        \item Bazı $x,N$ için $P=\lambd[x][N]$ ise bazı $N'$ için
          $P'=\lambd[x][N']$ ve $N\bredpar N'$ olmak zorundadır. Tümevarım
          varsayımına göre $N'\bredpar\bcd{N}$ ve $Q'\bredpar\bcd{Q}$.
          Dolayısıyla \olref{defn:bredpar4} gereği
          $(\lambd[x][N'])Q'\bredpar\Subst{\bcd{N}}{\bcd{Q}}{x}$.
        \item $P$ bir $\lambd$-soyutlaması değilse
          \olref{defn:bredpar3} gereği $P'Q'\bredpar\bcd{P}\bcd{Q}$;
          sağ taraf da \olref{defn:bcd3} gereği $\bcd{PQ}$'dur.
      \end{enumerate}
    \item Son kural \olref{defn:bredpar4} olsun. Bazı $x,N,Q,N',Q'$ için
      $M=(\lambd[x][N])Q$, $M'=\Subst{N'}{Q'}{x}$,
      $N\bredpar N'$ ve $Q\bredpar Q'$ olur. Tümevarım varsayımına göre
      $N'\bredpar\bcd{N}$ ve $Q'\bredpar\bcd{Q}$. \olref{lem:comp} gereği
      $\Subst{N'}{Q'}{x}\bredpar\Subst{\bcd{N}}{\bcd{Q}}{x}$; sağ taraf
      tam olarak $\bcd{((\lambd[x][N])Q)}$'dur.
  \end{enumerate}
\end{proof}

\begin{prob}
  \olref[lam][cr][pb]{lem:cont} lemasının ispatını tamamlayın.
\end{prob}

\begin{thm}\ollabel{thm:cr}
  $\bredpar$ Church--Rosser özelliğini sağlar.
\end{thm}

\begin{proof}
  Doğrudan \olref{lem:cont} lemasından çıkar.
\end{proof}

